\section{Provenance, endogenization, and routed realizations}\label{sec:provenance-routing}

The receiver interface specifies what reaches a continuation. A second question asks where the maps that construct that interface obtain their operative values. The answer requires semantic argument types, but it does not require a second ontology of state-indexed presentations. The actual routed law is given by the evaluated maps that refine $Q_j$.

\subsection{Marked argument origins and derived dependence}

Let an architectural component have a typed signature
\begin{equation}\label{eq:typed-component-signature}
c:
D_B\times D_\xi\times D_\theta\times X
\longrightarrow C.
\end{equation}
The arguments have declared meanings:
\[
\begin{aligned}
B&:\text{schema-declared structure},
&\xi&:\text{instance-supplied structure},\\
\theta&:\text{parameter data},
&H\in X&:\text{current represented state}.
\end{aligned}
\]
These semantic source types are marked. Actual dependence is derived relative to an admissible domain.

\begin{definition}[Essential dependence]\label{def:essential-dependence}
Let
\[
D\subseteq D_B\times D_\xi\times D_\theta\times X
\]
be an admissible domain. For $a\in\{B,\xi,\theta,H\}$, write
\[
a\in\Dep_D(c)
\]
when there exist $u,v\in D$ that agree in every coordinate except possibly $a$ and satisfy $c(u)\neq c(v)$.
\end{definition}

The definition distinguishes three claims:
\[
\boxed{
\begin{aligned}
\text{the schema permits an }H\text{-argument}
&\neq H\in\Dep_D(c),\\
H\in\Dep_D(c)
&\neq H\in\Dep_K(c)
\quad\text{for every observed }K\subseteq D.
\end{aligned}}
\]
A trained constructor can be state-dependent on its admissible domain while remaining constant on the states observed in one dataset or deployment window.

\begin{proposition}[Observed behavior underdetermines state provenance]\label{prop:provenance-underdetermination}
Let $\varnothing\neq K\subsetneq X$ and let $A_0\neq A_1$. Define
\[
\Phi_0(H)=A_0,
\qquad
\Phi_1(H)=
\begin{cases}
A_0,&H\in K,\\
A_1,&H\notin K.
\end{cases}
\]
Then $H\notin\Dep_X(\Phi_0)$ and $H\in\Dep_X(\Phi_1)$, while $\Phi_0=\Phi_1$ on $K$.
\end{proposition}

\begin{proof}
The first map is constant. The second takes two values because both $K$ and its complement are nonempty. Their agreement on $K$ is immediate.
\end{proof}

The extensional routing law on a restricted support therefore cannot establish whether its operative relation was supplied, learned but state-independent, or generated from state and accidentally constant there. The argument origins must be marked, and realized dependence must be tested on a declared admissible domain.

\subsection{Restriction and endogenization}

\begin{definition}[Restriction]\label{def:restriction-move}
A restriction replaces an admissible family $\Sigma$ for one typed architectural component by a subfamily
\[
\Sigma'\subseteq\Sigma.
\]
\end{definition}

\begin{definition}[Endogenization]\label{def:endogenization-move}
An endogenization replaces a state-independent component
\[
c(B,\xi,\theta)
\]
by a same-typed constructor
\[
\widetilde c(B,\xi,\theta,H).
\]
\end{definition}

Restriction changes which values or maps are admissible. Endogenization changes which arguments may construct the operative value. They are orthogonal. A fixed graph can carry state-dependent weights; a state-generated graph can still be restricted to bounded degree; a state-conditioned coefficient can remain constant on the realized domain.

The standard Transformer illustrates both moves. Relative to a supplied relation, it endogenizes receiver-relative allocation by making it a function of the current token states. It simultaneously restricts the construction to pairwise scalar evidence, positive support, receiver-row normalization, source-local values, additive aggregation, and a shared finite-rank score family.

\subsection{Operator-valued pair contributions}

Assume now that the source-resolved exposure of receiver $r$ is additive and vector-valued. Let $R$ be a finite receiver set and $S$ a finite source set. For each $i\in S$, let $U_i$ be a real vector space, and for each $r\in R$, let $W_r$ be a real receiver-exposure vector space.

\begin{definition}[Additive routed realization]\label{def:additive-routed-realization}
An additive routed realization of $E_r$ consists of payload maps
\[
u_i:X\to U_i
\]
and pair-contribution operators
\[
C_{ri}:X\to\mathcal L(U_i,W_r)
\]
such that
\begin{equation}\label{eq:additive-route}
\boxed{
E_r(H)
=
\sum_{i\in S}C_{ri}(H)u_i(H).
}
\end{equation}
The corresponding source-resolved presentation is
\[
P_r(H)=\bigl(C_{ri}(H)u_i(H)\bigr)_{i\in S},
\]
and $\Agg_r$ is summation.
\end{definition}

The operator $C_{ri}(H)$ is the minimal pair object in the additive linear regime. A further factorization
\begin{equation}\label{eq:scalar-transport-factorization}
C_{ri}(H)=\lambda_{ri}(H)T_{ri}(H)
\end{equation}
becomes architecturally meaningful only when scalar relevance and payload transport are separately marked components. Without that mark, the rescaling
\[
\lambda_{ri}\mapsto a_{ri}\lambda_{ri},
\qquad
T_{ri}\mapsto a_{ri}^{-1}T_{ri}
\]
leaves the pair operator unchanged wherever $a_{ri}\neq0$.

For fixed $H$, define
\begin{equation}\label{eq:state-operator}
A_H:\bigoplus_{i\in S}U_i\to\bigoplus_{r\in R}W_r,
\qquad
(A_Hu)_r:=\sum_iC_{ri}(H)u_i.
\end{equation}

\begin{proposition}[Routed interfaces define state-indexed operator families]\label{prop:routed-operator-field}
For every fixed $H$, $A_H$ is linear and
\[
E(H)=A_Hu(H).
\]
Thus the routed realization determines a state-to-operator map $H\mapsto A_H$.
\end{proposition}

\begin{proof}
Linearity follows from the linearity of every $C_{ri}(H)$ and of finite summation. Substituting the state-derived payload family gives \eqref{eq:additive-route} receiver by receiver.
\end{proof}

The current state performs two distinct roles: it supplies the payload $u(H)$ and may also construct the operator $A_H$ that acts on that payload. This state-to-operator lift is broader than attention.

\subsection{Support, contribution, and retained exposure}

Define the operative pair support
\[
\mathsf{Supp}_H:=\{(r,i):C_{ri}(H)\neq0\}.
\]
The following distinctions are structural:
\[
\boxed{
\begin{gathered}
\text{source carrier}
\neq
\text{admissible pair relation}
\neq
\text{operative support},\\
\text{pair operator}
\neq
\text{source payload}
\neq
\text{retained aggregate exposure}.
\end{gathered}}
\]
A nonzero pair operator may act on a zero payload. Several nonzero source contributions may cancel under summation. Different pair families may induce the same aggregate. These are reasons not to identify a routing matrix with direct mediation or with the complete information available to a receiver.

\subsection{Controlled contrasts}

The factorization coordinates place several familiar architectures without treating them as historical steps toward the Transformer.

\paragraph{Linear prediction.}
A linear predictor $x\mapsto Wx+b$ acts after a representation has already fixed $x$. At the one-receiver grain, the interface may be all of $x$ and the continuation is affine. The model does not by itself construct a receiver-relative source organization.

\paragraph{Normalized kernel estimation.}
A normalized kernel estimator
\[
\widehat y(x)
=
\sum_i
\frac{\kappa(x,x_i)}{\sum_j\kappa(x,x_j)}y_i
\]
already has a one-receiver source-resolved additive interface. Its allocation is query-conditioned, but the comparison geometry is normally supplied or globally parameterized and the result is usually one terminal observable rather than a successor representation that constructs a new routed law \citep{nadaraya1964estimating,watson1964smooth}.

\paragraph{Convolution.}
For a finite offset family $D$,
\[
(\operatorname{Conv}H)_r
=
\sum_{\delta\in D}W_\delta h_{r+\delta}.
\]
Here $C_{r,r+\delta}=W_\delta$ on a schema-fixed translation-generated support and zero elsewhere. The relation and parameter sharing are supplied by the schema; the payloads vary with the state.

\paragraph{Graph message passing.}
A graph layer
\[
E_r(H,\xi)
=
\Agg_{i\in N_\xi(r)}m_\theta(h_r,h_i,e_{ri})
\]
uses an instance-supplied admissible relation while allowing state-dependent messages or weights. Graph attention endogenizes allocation inside the supplied accessible sector rather than eliminating the graph \citep{velickovic2018graph,bronstein2021geometric}.

\paragraph{Recurrence and state-space summaries.}
A recurrent architecture gives the receiver a persistent summary instead of a source-resolved row over all predecessors. Long-range influence is mediated through that summary, so the interface has different fibers even when the external function can be matched.

\paragraph{Mixture-of-experts routing.}
A token state can construct coefficients over a supplied expert carrier. This is a routed interface, but its source and receiver types differ from token self-attention and its capacity constraints may couple allocations across tokens.

No scalar measure globally orders these systems. One may endogenize support while fixing aggregation, enlarge message types while compressing source resolution, or vary carrier populations while retaining a rigid schedule. The next section selects a narrower additive regime and derives the additional commitments that make it attention.
